% ch-a2-04 —— 第 2 篇 Zhu，PDF 91–103，印刷页 78–90
\ysrc{78}
……可能不为零）。其同构类的集合是 $B(Z)=H^{1}(W_F,Z(\overline F))$。为研究对偶一侧，把 $Z$ 嵌入一个 $F$-环面 $T$，并令 $T'=T/Z$。于是定义
\[
\hat Z := \hat T'/\hat T.
\]
若 $Z$ 是环面，则 $\hat Z$ 恰是 $Z$ 的对偶群；但一般情形下它只是一个 Picard 叠。例如若 $Z$ 有限，则 $\hat Z$ 是 $\ker(\hat T\to\hat T')$ 的分类叠。无论如何，$\hat Z$ 典范地不依赖于嵌入 $Z\to T$ 的选取，可称为 $Z$ 的对偶。

$\hat Z$ 上有 $\mathbb G_m\times\Gamma_{\wt F/F}$ 的自然作用（当然 $\mathbb G_m$ 的作用平凡，保留它是为了使记号统一）。于是可以定义
\[
{}^{c}Z := \hat Z\rtimes(\mathbb G_m\times\Gamma_{\wt F/F}),
\]
视之为 $\Z[1/p]$ 上的幺半叠。然后可以定义 $\Loc_{{}^{c}Z,F}$。这是一个 Picard 2-叠。也可以取其对偶 $\Loc^{\vee}_{{}^{c}Z,F}=\underline{\Hom}(\Loc_{{}^{c}Z,F},B\mathbb G_m)$。于是猜想 3.18 可作如下推广。
\begin{conj}{3.20}
导出 Picard 叠之间存在自然同构 $\mathbf{Tor}_{Z,\mathrm{iso}F}\cong\Loc^{\vee}_{{}^{c}Z,F}$。特别地，每个 $\theta\in\mathbf{Tor}_{Z,\mathrm{iso}F}$ 都给出 $\Loc_{{}^{c}Z,F}$ 上的一个乘法线丛 $\mathcal L_{\theta}$。
\end{conj}

把上述构造用于连通约化群 $G$ 的中心 $Z=Z_G$，以讨论 $\Coh(\Loc_{{}^{c}G,F,\iota})$ 的某种对称性。设 $\hat G_{sc}$ 是 $\hat G$ 的导群的单连通覆盖（即 $G_{ad}$ 的对偶群）。设 $\hat T_{sc}$ 是 $\hat T$ 在 $\hat G_{sc}$ 中的原像。则 $\hat Z_G\cong\hat T/\hat T_{sc}\cong\hat G/\hat G_{sc}$，因此有「行列式」映射 $\hat G\to\hat Z_G$，它诱导
\[
\delta:\Loc_{{}^{c}G,F,\iota}\to\Loc_{{}^{c}Z_G,F}.
\]
猜想 3.20 蕴含 $\mathbf{Tor}_{Z_G,\mathrm{iso}F}$ 在 $\Coh(\Loc_{{}^{c}G,F,\iota})$ 上有自然作用，由下式给出：
\begin{equation}
\mathbf{Tor}_{Z_G,\mathrm{iso}F}\times\Coh(\Loc_{{}^{c}G,F,\iota})\to\Coh(\Loc_{{}^{c}G,F,\iota}),\qquad (\theta,\mathcal F)\mapsto\delta^{*}\mathcal L_{\theta}\otimes\mathcal F. \tag{3.7}
\end{equation}
这是几何 Langlands 中某些构造的算术类比（例如见 [15, 3.8, 5.6]）。

可以把这个作用稍加细化。把 $Z_G\subset T$ 嵌入，得到一个映射
\[
B(Z_G)\to B(T)\cong\mathbb X^{\bullet}(\hat T^{\Gamma_F})\to\mathbb X^{\bullet}(Z_{\hat G}^{\Gamma_F}).
\]
复合映射 $B(Z_G)\to\mathbb X^{\bullet}(Z_{\hat G}^{\Gamma_F})$，$\theta\mapsto[\theta]$ 不依赖于 $T$ 的选取。另一方面，$\Loc_{{}^{c}G,F,\iota}$ 是一个 $Z_{\hat G}^{\Gamma_F}$-带叠（gerbe）（因为 $Z_{\hat G}^{\Gamma_F}\subset\hat G$ 平凡地作用在 $\Loc^{\square}_{{}^{c}G,F,\iota}$ 上）。由此可见存在分解
\begin{equation}
\Coh(\Loc_{{}^{c}G,F,\iota})=\bigoplus_{\beta\in\mathbb X^{\bullet}(Z_{\hat G}^{\Gamma_F})}\Coh^{\beta}(\Loc_{{}^{c}G,F,\iota}). \tag{3.8}
\end{equation}
于是作用 $\mathcal L_{\theta}$ 把 $\Coh^{\beta}(\Loc_{{}^{c}G,F,\iota})$ 送到 $\Coh^{\beta+[\theta]}(\Loc_{{}^{c}G,F,\iota})$。

$\Coh(\Loc_{{}^{c}G,F,\iota})$ 上还有另一种对称性。设 $\tau\in\Aut(\hat G,\hat B,\hat T,\hat e)$ 是 Cartan 对合，即使下式成立的唯一自同构：
\[
\tau:\mathbb X^{\bullet}(\hat T)\to\mathbb X^{\bullet}(\hat T),\qquad \lambda\mapsto\lambda^{*}=-w_0(\lambda),
\]
其中 $w_0$ 是 $\hat G$ 的 Weyl 群中长度最长的元素。由于 $\tau$ 在 $\Aut(\hat G,\hat B,\hat T,\hat e)$ 中中心，它诱导 ${}^cG$ 的一个自同构，从而诱导 $\Coh(\Loc_{{}^{c}G,F,\iota})$ 的一个自等价，仍记作 $\tau$。令
\begin{equation}
{}'\mathbb D^{\mathrm{Se}} := \tau\circ\mathbb D^{\mathrm{Se}}:\Coh(\Loc_{{}^{c}G,F,\iota})\to\Coh(\Loc_{{}^{c}G,F,\iota}). \tag{3.9}
\end{equation}

\ysrc{79}
为修正的 Grothendieck–Serre 对偶。注意 ${}'\mathbb D^{\mathrm{Se}}$ 保持分解 (3.8) 并与作用 (3.7) 交换，而原来的 Grothendieck–Serre 对偶函子 $\mathbb D^{\mathrm{Se}}:\Coh(\Loc_{{}^{c}G,F,\iota})\to\Coh(\Loc_{{}^{c}G,F,\iota})$ 则不然。

\subsection{谱抛物诱导}

设 $\hat P$ 是 $\hat G$ 的抛物子群，包含 $\hat B$，且在 $\Gamma_{\wt F/F}$ 于 $\hat G$ 上的作用下稳定；设 $\hat M$ 是其标准 Levi 因子（即包含 $\hat T$ 的那个）。则 $\mathbb G_m\times\Gamma_{\wt F/F}$ 在 $\hat G$ 上的作用保持 $\hat P$ 与 $\hat M$，因此可分别构造 ${}^cP$ 与 ${}^cM$，并类似地定义 $\Loc_{{}^{c}P,F,\iota}$ 与 $\Loc_{{}^{c}M,F,\iota}$。注意，与 $\Loc_{{}^{c}G,F,\iota}$、$\Loc_{{}^{c}M,F,\iota}$ 不同，$\Loc_{{}^{c}P,F,\iota}$ 可能不是经典的（见注 2.34），尽管它仍是拟光滑的。\tn{底本此处作「may not be not classical」，多了一重否定，疑为「may not be classical」之误。}我们强调，在下面的讨论中必须记住 $\Loc_{{}^{c}P,F,\iota}$ 的导出结构。$\Z[1/p]$ 上有如下交换图：
\begin{equation}
\begin{tikzcd}[column sep=large,row sep=large]
& \Loc_{{}^{c}P,F,\iota} \arrow[dl, "r"'] \arrow[dr, "\pi"] & \\
\Loc_{{}^{c}M,F,\iota} \arrow[ur, "i"'] \arrow[d] & & \Loc_{{}^{c}G,F,\iota} \arrow[d] \\
\Spec Z_{{}^{c}M,F} \arrow[rr] & & \Spec Z_{{}^{c}G,F}
\end{tikzcd} \tag{3.10}
\end{equation}
其中 $\pi,r,i$ 由 $\hat G,\hat P,\hat M$ 之间相应的态射诱导，底部的映射由 $\pi\circ i:\Loc_{{}^{c}M,F,\iota}\to\Loc_{{}^{c}G,F,\iota}$ 诱导。为看出此图交换，只需说明 $r$ 诱导同构
\begin{equation}
H^{0\Gamma}(\Loc_{{}^{c}M,F,\iota},\mathcal O)\to H^{0\Gamma}({}^{cl}\Loc_{{}^{c}P,F,\iota},\mathcal O). \tag{3.11}
\end{equation}
令 $2\rho_{\hat G,\hat M}=2\rho-2\rho_{\hat M}$，其中 $2\rho$（相应地 $2\rho_{\hat M}$）是 $\hat G$（相应地 $\hat M$）的正余根之和。则 $2\rho_{\hat G,\hat M}(\mathbb G_m)$ 在 ${}^cP$ 上的共轭作用把它收缩到 ${}^cM$ 中。等价地，$\Lambda[{}^cP]$ 关于 $2\rho_{\hat G,\hat M}(\mathbb G_m)$ 的零权部分恰是 $\Lambda[{}^cM]$。由此可见 (3.11) 是同构。

若令 $W_{{}^{c}G,{}^{c}M}$ 是 ${}^cM\subset{}^cG$ 在 $\hat G$ 中的正规化子对 $\hat M$ 的商，则映射 $Z_{{}^{c}G,F}\to Z_{{}^{c}M,F}$ 通过
\begin{equation}
Z_{{}^{c}G,F}\to(Z_{{}^{c}M,F})^{W_{{}^{c}G,{}^{c}M}} \tag{3.12}
\end{equation}
分解。

我们有下述引理（比较 [2, 13.2.2]）。
\begin{lem}{3.21}
态射 $r$ 拟光滑，$\pi$ 本征（proper）且概形（schematic）。
\end{lem}
\begin{proof}
$\pi$ 本征且概形是清楚的。至于 $r$ 的拟光滑性，只需注意 $\rho\in\Loc_{{}^{c}P,F,\iota}$ 处的相对余切复形是 $C_{*}(\Gamma_{F,\iota},\mathrm{Ad}^{u,*}_{\rho})[-1]$；当 $\rho$ 是经典点时，它集中在次数 $[-1,1]$。这里 $\mathrm{Ad}^{u,*}$ 是 ${}^cP$ 在其幂幺根的李代数的对偶上的余伴随表示。
\end{proof}

回忆 Arinkin–Gaitsgory（在 [2] 中）对特征零域上的拟光滑导出代数叠 $X$ 附以一个经典奇点叠 $\Sing(X)$，并对 $X$ 上的凝聚层 $\mathcal F$ 附以一个锥子集 $\Sing(\mathcal F)\subset\Sing(X)$ 作为它的奇点支撑。这些构造对

\ysrc{80}
$\mathbf{CAlg}_{\Lambda}$ 上的叠稍作修改后仍然成立（细节见 [88, §9.4]）。特别地，按定义
\[
\Sing(\Loc_{{}^{c}G,F,\iota})=\{(\rho,\xi)\mid \rho\in{}^{cl}\Loc_{{}^{c}G,F,\iota},\ \xi\in H_{2}(\Gamma_{F,\iota},\mathrm{Ad}^{*}_{\rho})\},
\]
其中 $\mathrm{Ad}^{*}$ 表示 ${}^cG$ 在 $\hat G$ 的李代数的对偶上的余伴随表示。

如 [2] 所解释，一个特殊的锥子集 $\hat{\mathcal N}_{{}^{c}G,F,\iota}\subset\Sing(\Loc_{{}^{c}G,F,\iota})$ 在 Langlands 对应中起重要作用。利用 (2.24)（或局部 Tate 对偶的某个版本），有
\[
H_{2}(\Gamma_{F,\iota},\mathrm{Ad}^{*}_{\rho})\cong(\hat{\mathfrak g}^{*})^{\rho(I_{F,\iota})=1,\ \rho(\sigma)=q^{-1}}\subset\mathrm{Ad}^{*}_{\rho}.
\]
设 $\hat{\mathcal N}^{*}\subset\hat{\mathfrak g}^{*}$ 是 $\hat{\mathfrak g}^{*}$ 的幂零锥。定义
\begin{equation}
\hat{\mathcal N}_{{}^{c}G,F,\iota}=\{(\rho,\xi)\in\Sing(\Loc_{{}^{c}G,F,\iota}),\ \xi\in\hat{\mathcal N}^{*}_{\rho}\}. \tag{3.13}
\end{equation}

下述命题可完全仿照 [2, 13.2.6] 证明。回忆我们对 $\Loc_{{}^{c}G,F,\iota}$ 上凝聚层的约定（见注 3.8）。
\begin{prop}{3.22}
存在良定义的函子（称为谱抛物诱导）
\[
\pi_{*}r^{!}:\Coh(\Loc_{{}^{c}M,F,\iota})\to\Coh(\Loc_{{}^{c}G,F,\iota}),
\]
它限制为函子 $\pi_{*}r^{!}:\Coh_{\hat{\mathcal N}_{{}^{c}M,F,\iota}}(\Loc_{{}^{c}M,F,\iota})\to\Coh_{\hat{\mathcal N}_{{}^{c}G,F,\iota}}(\Loc_{{}^{c}G,F,\iota})$。
\end{prop}

我们有如下观察。
\begin{lem}{3.23}
在 $\Q$ 上，$\Sing(\Loc_{{}^{c}G,F,\iota}\otimes\Q)=\hat{\mathcal N}_{{}^{c}G,F,\iota}\otimes\Q$。然而在 $\mathbb F_{\ell}$ 上，当 $\ell\mid q-1$ 时，$\Sing(\Loc_{{}^{c}G,F,\iota})$ 严格大于 $\hat{\mathcal N}_{{}^{c}G,F,\iota}$。
\end{lem}
\begin{proof}
利用引理 3.10 中 $\Loc_{{}^{c}G,F,\iota}\otimes\Q$ 与 $\Loc^{WD}_{{}^{c}G,F}$ 之间的等同，把 $H_{2}(\Gamma_{F,\iota},\mathrm{Ad}^{*}_{\rho})$ 等同于
\[
\{\xi\in(\hat{\mathfrak g}^{*})^{r(I_F)}\mid \mathrm{ad}^{*}_{X}(\xi)=0,\ r(\sigma)(\xi)=q^{-1}\xi\},
\]
其中 $(r,X)$ 按引理 3.10 对应于 $\rho$。需要说明这样的 $\xi$ 自动幂零。令 $\mathfrak h:=\hat{\mathfrak g}^{r(I_F)}$，它是约化李代数。可以把 $(\hat{\mathfrak g}^{*})^{r(I_F)}$ 与 $\mathfrak h$ 等同，视作 $(r(\sigma),\mathfrak h)$-模。于是 $\mathrm{ad}^{*}_{\xi}(\xi)$ 是 $r(\sigma)$ 的特征向量，特征值为 $q^{-j-1}$。这迫使 $\mathrm{ad}^{j}_{\xi}(\xi)=0$，只要 $j$ 足够大。也就是说，$\xi$ 幂零。
\end{proof}

上述计算还蕴含以下结论。
\begin{lem}{3.24}
设 $\rho:W_F\to{}^cG(\overline\Q_{\ell})$ 是连续表示，且 $\mathrm{Ad}^{0}_{\rho}:W_F\to\mathrm{GL}(\hat{\mathfrak g})$ 是权零纯的（Deligne 意义下），则 $\rho$ 是 $\Loc_{{}^{c}G,F}$ 的光滑点。
\end{lem}
\begin{proof}
事实上，此时 $H^{2}(W_F,\mathrm{Ad}^{0}_{\rho})=0$，应用命题 2.19 即可得证。
\end{proof}

在本小节的剩余部分，我们假定 $\wt F/F$ 驯顺分歧，即 $P_F\subset\Gamma_F\to\Gamma_{\wt F/F}$ 的像是平凡的。于是有叠 $\Loc_{{}^{c}G,F^{t}/F,\iota}$，称为驯顺 Langlands 参数的叠，也记作 $\Loc^{\mathrm{tame}}_{{}^{c}G,F,\iota}$。它是 $\Loc_{{}^{c}G,F,\iota}$ 的开闭子叠。

\ysrc{81}
设 $\Loc^{\mathrm{tame},\square}_{{}^{c}G,F,\iota}$ 记加框架的版本。具体地说，若把 $\tau$（相应地 $\sigma$）在映射 $\Gamma_q\xrightarrow{\iota}\Gamma^{t}_F\to\Gamma_{\wt F/F}$ 下的像记作 $\bar\tau$（相应地 $\bar\sigma$），则
\begin{equation}
\Loc^{\mathrm{tame},\square}_{{}^{c}G,F,\iota}\cong\{(\tau,\sigma)\in\hat G^{\bar\tau}\times\hat G^{q^{-1}}\bar\sigma\mid \sigma\tau\sigma^{-1}=\tau^{q}\}\subset{}^cG\times{}^cG. \tag{3.14}
\end{equation}
\begin{rem}{3.25}
可以把 $\Loc^{\mathrm{tame},\square}_{{}^{c}G,F,\iota}$ 与 $\hat G$ 的交换概形相比较，后者分类 $\hat G$ 中彼此交换的元素对。虽然这两个叠在 $\Q$ 上呈现出很不同的几何结构，但在 $\mathbb F_{\ell}$ 上、当 $\ell\mid q-1$ 时它们具有一些相似的性质。
\end{rem}

可以类似地定义 $\Loc^{\mathrm{tame}}_{{}^{c}B,F,\iota}$ 与 $\Loc^{\mathrm{tame}}_{{}^{c}T,F,\iota}$。注意 $\Loc^{\mathrm{tame}}_{{}^{c}T,F,\iota}$ 恰是例 3.17 中讨论的 $\Loc^{(1)}_{{}^{c}T,F}$。存在与 (3.10) 类似的图，只是各处都添加上标 $(-)^{\mathrm{tame}}$。与引理 3.21 相同的论证表明 $r^{\mathrm{tame}}$ 拟光滑，$\pi^{\mathrm{tame}}$ 本征且概形。

包含 $\langle\tau\rangle\subset\Gamma_q$ 诱导映射
\begin{equation}
\Loc^{\mathrm{tame}}_{{}^{c}G,F,\iota}\to\hat G^{\bar\tau}/\hat G\to\hat G^{\bar\tau}\sslash\hat G\cong\hat A\sslash W_0, \tag{3.15}
\end{equation}
其中 $\hat A=\hat T/(1-\bar\tau)\hat T$，$W_0=W^{\bar\tau}$ 是 $\hat G$ 的 Weyl 群 $W$ 的 $\bar\tau$-不变量（例如见 [78, 4.2.3]）。第二个映射是 GIT 商映射，最后的同构是 Chevalley 限制同构。如命题 2.33 的证明所示，该态射通过 $\Loc^{\mathrm{tame}}_{{}^{c}G,F,\iota}\to(\hat A\sslash W_0)^{[q]}$ 分解，其中 $(\hat A\sslash W_0)^{[q]}$ 是映射 $[q]:\hat A\sslash W_0\to\hat A\sslash W_0$ 的（经典）不动点子概形，而 $[q]$ 由态射 $\hat G^{\bar\tau}\to\hat G^{\bar\tau}$，$g\bar\tau\mapsto\bar\sigma^{-1}(g\bar\tau)^{q}\bar\sigma$ 诱导。不难验证 $(\hat A\sslash W_0)^{[q]}$ 在 $\Z[1/p]$ 上有限，且在 $\Q$ 上 étale。令 $1:\Spec\Z[1/p]\to\hat A\sslash W_0$ 是对应于 $\hat A$ 的单位元的映射，$\{1\}^{\wedge}$ 记 $\hat A\sslash W_0$ 沿 $\{1\}$ 的形式完备化。

我们定义幂幺参数叠的两个版本：
\[
\Loc^{\mathrm{unip}}_{{}^{c}G,F,\iota}:=\Loc^{\mathrm{tame}}_{{}^{c}G,F,\iota}\times_{\hat A\sslash W_0}\{1\}\subset\Loc^{\widehat{\mathrm{unip}}}_{{}^{c}G,F,\iota}:=\Loc^{\mathrm{tame}}_{{}^{c}G,F,\iota}\times_{\hat A\sslash W_0}\{1\}^{\wedge}.
\]
\begin{rem}{3.26}
(1) 按定义，$\Loc^{\mathrm{unip}}_{{}^{c}G,F,\iota}$ 是代数叠，但一般是导出的。另一方面，$\Loc^{\widehat{\mathrm{unip}}}_{{}^{c}G,F,\iota}$ 是经典的，但为 ind-代数叠。二者底层的约化子叠一致：
\[
{}^{\mathrm{red}}\Loc^{\mathrm{unip}}_{{}^{c}G,F,\iota}={}^{\mathrm{red}}\Loc^{\widehat{\mathrm{unip}}}_{{}^{c}G,F,\iota}.
\]
底变到一个域时，$\Loc^{\widehat{\mathrm{unip}}}_{{}^{c}G,F,\iota}$ 实际上是代数叠（见 [88, Lemma 2.14]），但一般不约化。事实上连 ${}^{cl}\Loc^{\mathrm{unip}}_{{}^{c}G,F,\iota}$ 都可能不约化（例如 $G=\mathrm{PGL}_2$ 的情形见 [26]）。

在 $\Q$ 上情形要好得多。由于 $1$ 是 $(\hat A\sslash W_0)^{[q]}\otimes\Q$ 的孤立点，可见
\[
\Loc^{\widehat{\mathrm{unip}}}_{{}^{c}G,F,\iota}\otimes\Q=(\Loc^{\mathrm{tame}}_{{}^{c}G,F,\iota}\times_{(\hat A\sslash W_0)^{[q]}}\{1\})\otimes\Q
\]
在 $\Loc^{\mathrm{tame}}_{{}^{c}G,F,\iota}\otimes\Q$ 中开闭。当 $\bar\tau=1$ 时，$\Loc^{\widehat{\mathrm{unip}}}_{{}^{c}G,F,\iota}\otimes\Q$ 是 $\Loc^{\mathrm{tame}}_{{}^{c}G,F,\iota}\otimes\Q$ 的一个连通分支。特别地，它仍是局部完全交。
\end{rem}

\ysrc{82}
(2) 我们的术语可能有些误导：对一个（取值于域的）点 $\rho\in\Loc^{\widehat{\mathrm{unip}}}_{{}^{c}G,F,\iota}$，元素 $\rho(\tau)\in\hat G^{\bar\tau}$ 未必是幂幺元素（因为 $\bar\tau$ 未必平凡）。另一方面，若 $\bar\tau=1$，即 $\wt F/F$ 非分歧，则
\[
\Loc^{\widehat{\mathrm{unip}}}_{{}^{c}G,F,\iota}\cong\Loc^{\widehat{\mathrm{unip}},\square}_{{}^{c}G,F,\iota}/\hat G,
\quad\text{其中}\quad
\Loc^{\widehat{\mathrm{unip}},\square}_{{}^{c}G,F}=\{(\tau,\sigma)\in\hat{\mathcal U}^{\wedge}\times\hat G^{q^{-1}}\bar\sigma\mid \sigma\tau\sigma^{-1}=\tau^{q}\},
\]
其中如前 $\hat{\mathcal U}$ 是 $\hat G$ 的幂幺簇，$\hat{\mathcal U}^{\wedge}$ 记它在 $\hat G$ 中的形式完备化。于是 $\tau$ 在 $\hat G$ 中的像确为幂幺的。

若 $\wt F/F$ 非分歧，则 $\Loc^{\mathrm{unip}}_{{}^{c}G,F,\iota}$ 内部有非分歧参数的叠
\[
\Loc^{\mathrm{ur},\square}_{{}^{c}G,F}\cong\hat G^{q^{-1}}\bar\sigma\subset{}^cG,\qquad \Loc^{\mathrm{ur}}_{{}^{c}G,F}=\Loc^{\mathrm{ur},\square}_{{}^{c}G,F}/\hat G.
\]
注意这个叠光滑，且不依赖于 $\iota$ 的选取（因此把它从记号中略去）。若 $T$ 是非分歧环面，则
\begin{equation}
\Loc^{\mathrm{ur}}_{{}^{c}T,F}={}^{\mathrm{red}}\Loc^{\mathrm{unip}}_{{}^{c}T,F}={}^{cl}\Loc^{\mathrm{unip}}_{{}^{c}T,F} \tag{3.16}
\end{equation}
与 §3.2 中讨论的 $\Loc^{(0)}_{{}^{c}T,F}$ 一致。

在本小节末尾，我们引入所谓谱 Deligne–Lusztig 叠及其幂幺版本。回忆我们假定 $\wt F/F$ 驯顺分歧。但建议读者初次阅读时先在 $\wt F/F$ 非分歧（从而 $\bar\tau=1$）的较简单情形下过一遍构造。

设 $\tilde{\hat G}^{\bar\tau}:=\hat G\times_{\hat B}\hat B^{\bar\tau}\to\hat G^{\bar\tau}$ 是 $\hat G^{\bar\tau}$ 的（扭）Grothendieck–Springer 消解（例如见 [77, 5.3]）。于是定义（大）Steinberg 簇 $\mathrm{St}_{\hat G^{\bar\tau}}=\tilde{\hat G}^{\bar\tau}\times_{\hat G^{\bar\tau}}\tilde{\hat G}^{\bar\tau}$，它是维数为 $\dim\hat G$ 的经典、约化、局部完全交概形。它的不可约分支自然地由 $W_0=W^{\bar\tau}$ 编号。对 $w\in W_0$，令 $\mathrm{St}_{\hat G^{\bar\tau},w}$ 记相应的不可约分支。为简单起见，记 $S=\mathrm{St}_{\hat G^{\bar\tau}}/\hat G$，$S_w=\mathrm{St}_{\hat G^{\bar\tau},w}/\hat G$。称 $S$ 为大 Steinberg 叠。

回忆 (3.15) 中的态射 $\Loc^{\mathrm{tame}}_{{}^{c}G,F,\iota}\to\hat G^{\bar\tau}/\hat G$。于是定义
\begin{equation}
\widetilde{\Loc}^{\mathrm{tame}}_{{}^{c}G,F,\iota}:=\Loc^{\mathrm{tame}}_{{}^{c}G,F,\iota}\times_{\hat G^{\bar\tau}/\hat G}\hat B^{\bar\tau}/\hat B\xrightarrow{\tilde\pi\times\mathrm{pr}}\Loc^{\mathrm{tame}}_{{}^{c}G,F,\iota}\times\hat B^{\bar\tau}/\hat B. \tag{3.17}
\end{equation}
于是 ${}^{cl}\widetilde{\Loc}^{\mathrm{tame}}_{{}^{c}G,F,\iota}$ 分类 $(\tau,\sigma,g\hat B)$，其中 $(\tau,\sigma)$ 如 (3.14) 是驯顺 Langlands 参数，$g\hat B\in\hat G/\hat B$ 使得 $\tau\in g^{-1}(\hat B^{\bar\tau})g$。注意由于 $\tau\in(g\sigma)^{-1}(\hat B^{\bar\tau})(g\sigma)$，还有另一个投影 $\mathrm{pr}':\widetilde{\Loc}^{\mathrm{tame}}_{{}^{c}G,F,\iota}\to\hat B^{\bar\tau}/\hat B$。因此有态射
\[
\widetilde{\Loc}^{\mathrm{tame}}_{{}^{c}G,F,\iota}\xrightarrow{\mathrm{pr}\times\mathrm{pr}'}\hat B^{\bar\tau}/\hat B\times_{\hat G^{\bar\tau}/\hat G}\hat B^{\bar\tau}/\hat B\cong S.
\]
于是定义
\begin{equation}
\widetilde{\Loc}^{\mathrm{tame},w}_{{}^{c}G,F,\iota}:=\widetilde{\Loc}^{\mathrm{tame}}_{{}^{c}G,F,\iota}\times_{S}S_w. \tag{3.18}
\end{equation}

\ysrc{83}
\begin{rem}{3.27}
若 $w=1$ 是单位元，可以证明
\[
\widetilde{\Loc}^{\mathrm{tame},1}_{{}^{c}G,F,\iota}\cong\Loc^{\mathrm{tame}}_{{}^{c}B,F,\iota}.
\]
非正式地说，$\widetilde{\Loc}^{\mathrm{tame},w}_{{}^{c}G,F,\iota}$ 分类那些 $(\tau,\sigma,B')$，使得 $B'$ 与 $\sigma B'\sigma^{-1}$ 的相对位置被 $w$ 界定。因此，一般的 $\widetilde{\Loc}^{\mathrm{tame},w}_{{}^{c}G,F,\iota}$ 可以称为谱 Deligne–Lusztig 叠。
\end{rem}

我们还引入谱 Deligne–Lusztig 叠的幂幺版本。考虑映射
\[
\widetilde{\Loc}^{\mathrm{tame}}_{{}^{c}G,F,\iota}\to\hat B^{\bar\tau}/\hat B\to\hat T^{\bar\tau}/\hat T\to\hat A=\hat T/(1-\bar\tau)\hat T.
\]
先定义
\[
(\hat B^{\bar\tau}/\hat B)^{\mathrm{unip}}=(\hat B^{\bar\tau}/\hat B)\times_{\hat A}\{1\}.
\]
注意若 $\bar\tau=1$，则 $(\hat B^{\bar\tau}/\hat B)^{\mathrm{unip}}=\hat U/\hat B$，其中 $\hat U$ 是 $\hat B$ 的幂幺根。然后定义 Steinberg 叠的幂幺版本
\begin{equation}
S^{\mathrm{unip}}=(\hat B^{\bar\tau}/\hat B)^{\mathrm{unip}}\times_{\hat G^{\bar\tau}/\hat G}(\hat B^{\bar\tau}/\hat B)^{\mathrm{unip}}, \tag{3.19}
\end{equation}
以及
\begin{equation}
\widetilde{\Loc}^{\mathrm{unip}}_{{}^{c}G,F,\iota}:=\widetilde{\Loc}^{\mathrm{tame}}_{{}^{c}G,F,\iota}\times_{\hat A}\{1\}. \tag{3.20}
\end{equation}
类似地有映射
\[
\widetilde{\Loc}^{\mathrm{unip}}_{{}^{c}G,F,\iota}\to S^{\mathrm{unip}}.
\]
对 $w\in W_0$，令
\[
S^{\mathrm{unip}}_w:=((\hat B^{\bar\tau}/\hat B)\times_{\hat G^{\bar\tau}/\hat G}(\hat B^{\bar\tau}/\hat B)^{\mathrm{unip}})\cap S_w,
\]
其中交集在 $S$ 中取。它是经典叠，虽然一般不可约。此外，容易看出映射
\[
S^{\mathrm{unip}}_w\subset(\hat B^{\bar\tau}/\hat B)\times_{\hat G^{\bar\tau}/\hat G}(\hat B^{\bar\tau}/\hat B)^{\mathrm{unip}}
\]
通过
\[
S^{\mathrm{unip}}_w\subset S^{\mathrm{unip}}\subset(\hat B^{\bar\tau}/\hat B)\times_{\hat G^{\bar\tau}/\hat G}(\hat B^{\bar\tau}/\hat B)^{\mathrm{unip}}
\]
分解。于是与 (3.18) 类似，可以定义
\begin{equation}
\widetilde{\Loc}^{\mathrm{unip},w}_{{}^{c}G,F,\iota}:=\widetilde{\Loc}^{\mathrm{unip}}_{{}^{c}G,F,\iota}\times_{S^{\mathrm{unip}}}S^{\mathrm{unip}}_w. \tag{3.21}
\end{equation}
与注 3.27 类似，可以证明
\[
\widetilde{\Loc}^{\mathrm{unip},1}_{{}^{c}G,F,\iota}\cong\Loc^{\mathrm{unip}}_{{}^{c}B,F,\iota},
\]
其中定义
\begin{equation}
\Loc^{\mathrm{unip}}_{{}^{c}B,F,\iota}:=\Loc^{\mathrm{tame}}_{{}^{c}B,F,\iota}\times_{\Loc^{\mathrm{tame}}_{{}^{c}T,F,\iota}}{}^{cl}\Loc^{\mathrm{unip}}_{{}^{c}T,F,\iota}. \tag{3.22}
\end{equation}
\ysrc{84}
\subsection{整体 Langlands 参数的叠}

现在转向整体 Langlands 参数。目前我们不知道如何在 $\Z$ 上（或在特征 $p$ 的函数域的 $\Z[1/p]$ 上）定义整体 Langlands 参数的叠，因此没有 $\Loc_{{}^{c}G,F,\iota}$ 的整体类比。不过本小节的主要目标是说明，第 2.4 节的一般配方在整体函数域情形下给出 $\Spec\Z_{\ell}$ 上一个合理的叠定义。数域情形更复杂，是与 Emerton 的进行中的合作 [26]；我们只在本小节末尾简要讨论。

固定一些记号。设 $F$ 是整体域。把伽罗瓦群 $\Gamma_F$ 视为投射有限群，而在整体函数域情形把 Weil 群 $W_F$ 视为局部投射有限群。令 $\Lambda=\Z_{\ell}$，其中若 $F$ 是函数域则 $\ell\neq\mathrm{char}\,F$。对素位 $v$，令 $F_v$ 记相应的局部域，$\kappa_v$ 记剩余域，$q_v=\sharp\kappa_v$。令 $\Gamma_v$（相应地 $W_v$）记 $F_v$ 的伽罗瓦（相应地 Weil）群。设 $G$ 是 $F$ 上的连通约化群。记 $G_v$ 为 $G_{F_v}$ 或 $G(F_v)$。$G$ 的 C-群记作 ${}^cG$，$G_v$ 的 C-群记作 ${}^cG_v$。对不位于 $\ell$ 之上的素位 $v$，简单起见把 $\Loc^{?}_{v}$ 记作 $\Loc^{?}_{{}^{c}G_v,F_v}$，其中 $?=\emptyset,\mathrm{tame},\mathrm{ur}$ 等。我们固定一个非空的有限素位集 $S$，包含所有无穷素位、$\ell$ 之上的素位以及在 $\wt F/F$ 中分歧的素位，并考虑对应于 $F$ 的在 $S$ 外非分歧的最大伽罗瓦扩张的商 $\Gamma_{F,S}$。类似地，在整体函数域情形有 $W_{F,S}$。设 $Y$ 是以 $F$ 为分式域、以 $\Gamma_{F,S}$ 为 étale 基本群的 Dedekind 概形。

现在设 $F$ 是函数域。令 $\mathbb F_q$ 是 $\mathbb F_p$ 在 $F$ 中的代数闭包。则 $Y$ 是 $\mathbb F_q$ 上的仿射光滑曲线。令 $\overline Y$ 是 $Y$ 到 $\overline{\mathbb F}_q$ 的底变。令 $\pi_1(\overline Y)$ 记几何基本群。（我们忽略 $Y$ 上基点的选取，因为它在下文作用不大。）回忆短正合列
\[
1\to\pi_1(\overline Y)\to W_{F,S}\xrightarrow{\|\cdot\|}\Z=\langle\sigma\rangle\to1.
\]
把 (3.2) 中的局部 Weil 群 $W_F$ 换成 $W_{F,S}$，定义
\begin{equation}
\Loc^{\wedge,\square}_{{}^{c}G,F,S}:=\mathcal R_{W_{F,S},{}^{c}G}\times_{\mathcal R_{W_{F,S},\mathbb G_m\times\Gamma_{\wt F/F}}}\{\chi\},\qquad
\Loc^{\wedge}_{{}^{c}G,F,S}=\Loc^{\wedge,\square}_{{}^{c}G,F,S}/\hat G^{\wedge}_{\ell}. \tag{3.23}
\end{equation}
令 $\Loc_{{}^{c}G,F,S,r}=\Loc^{\square}_{{}^{c}G,F,S,r}/\hat G_r$ 是 (3.23) 到 $\Spec\Z/\ell^{r}$ 的限制。则 $\Loc^{\square}_{{}^{c}G,F,S,r}$ 对每个 $\Z/\ell^{r}$-代数 $A$ 分类从 $W_{F,S}$ 到 ${}^cG(A)$ 的连续同态 $\rho$ 的空间，使得 $d\circ\rho=\chi$（引理 2.39）。还可以用定义 2.56 把 (3.23) 推广到 $\Spf\Z_{\ell}$：
\begin{equation}
\Loc^{\square}_{{}^{c}G,F,S}:=\mathcal R^{sc}_{W_{F,S},{}^{c}G}\times_{\mathcal R^{sc}_{W_{F,S},\mathbb G_m\times\Gamma_{\wt F/F}}}\{\chi\},\qquad
\Loc_{{}^{c}G,F,S}=\Loc^{\square}_{{}^{c}G,F,S}/\hat G. \tag{3.24}
\end{equation}
\begin{rem}{3.28}
最近 [3] 对函数域提出了 $\Q_{\ell}$ 上整体 Langlands 参数叠的另一种定义。他们的定义与上面给出的不同，但很可能给出一个与我们的 $\Loc_{{}^{c}G,F,S}$ 到 $\Q_{\ell}$ 的底变同构的叠。
\end{rem}

本小节的主结果如下。
\begin{thm}{3.29}
设 $\ell>2$。则 $\Loc_{{}^{c}G,F,S}$ 是 $\Z_{\ell}$ 上的拟光滑代数叠。它分解为开闭子叠的不交并
\begin{equation}
\Loc_{{}^{c}G,F,S}=\bigsqcup_{\Theta}\Loc^{\Theta}_{{}^{c}G,F,S}, \tag{3.25}
\end{equation}
\ysrc{85}
其中 $\Theta$ 取遍 $\mathcal R^{sc}_{\pi_1(\overline Y),{}^{c}G^{\bullet}}\sslash\hat G$ 中满足 $d\circ\Theta=\chi$ 的所有闭点。每个 $\Loc^{\Theta}_{{}^{c}G,F,S}$ 拟紧，且对 $\Loc^{\Theta}_{{}^{c}G,F,S}$ 的每个 $\mathbb F_{\ell}$-或 $\Q_{\ell}$-点 $x$，（剩余）伪表示 $\bar\rho_{x}|_{\pi_1(\overline Y)}$ 就是 $\Theta$。
\end{thm}

记号 $\bar\rho_{x}|_{\pi_1(\overline Y)}$ 见引理 2.60 及其前的讨论。

为证明该定理，先回忆 de Jong 猜想（[22]）：若 $\rho:\pi_1(\overline Y)\to\mathrm{GL}_m(\kappa((t)))$ 是算术基本群的连续表示，其中 $\kappa$ 是特征 $\ell$ 的有限域、$\kappa((t))$ 带 $t$-进拓扑，则 $\rho(\pi_1(\overline Y))$ 有限。这在 $\ell>2$ 的假设下由 Gaitsgory [30] 证明（另见 [8]）。\footnote{这也是我们要求 $\ell>2$ 的原因。这一限制理应可以去掉。}注意在 de Jong 猜想的陈述中可以把 $\pi_1(\overline Y)$ 换成 Weil 群 $W_{F,S}$。

我们需要下述推论。由于 Frobenius $\sigma$ 通过（外）自同构作用在 $\pi_1(\overline Y)$ 上，它也作用在 $\pi_1(\overline Y)$ 的伪表示的空间 $\mathcal R^{c}_{\pi_1(\overline Y),\mathrm{GL}^{\bullet}_m}\sslash\mathrm{GL}_m$ 上。设 $\mathcal R^{c,\Theta}_{\pi_1(\overline Y),\mathrm{GL}^{\bullet}_m}\sslash\mathrm{GL}_m$ 是一个 $\sigma$-稳定的连通分支。回忆它是导出形式概形。把其下经典形式概形的函数环写作
\[
A^{\Theta}:=\Gamma({}^{cl}\mathcal R^{c,\Theta}_{\pi_1(\overline Y),\mathrm{GL}^{\bullet}_m}\sslash\mathrm{GL}_m,\mathcal O).
\]
由于 $\pi_1(\overline Y)$ 满足 Mazur 条件 $\Phi_{\ell}$，这是一个完备诺特局部 $\Z_{\ell}$-代数（[16, 3.7]），$\sigma$ 在其上有作用。
\begin{lem}{3.30}
商环 $A^{\Theta}/(\sigma-1)A^{\Theta}$ 在 $\Z_{\ell}$ 上有限。
\end{lem}
\begin{proof}
注意 $B^{\Theta}=A^{\Theta}/(\sigma-1)A^{\Theta}$ 仍是剩余域为 $\kappa$ 的完备诺特局部环。因此只需证明 $B^{\Theta}/\ell$ 是阿廷环。设 $B^{\Theta}\to\kappa'[[t]]$ 是局部环同态，$\kappa'$ 在 $\kappa$ 上有限，它给出 $\pi_1(\overline Y)$ 的一个连续 $\kappa'[[t]]$-值伪表示。则由命题 2.49，这样的 $\kappa'((t))$-值伪表示来自某个有限扩张 $K/\kappa'((t))$ 上的连续（绝对）半单表示 $\rho:\pi_1(\overline Y)\to\mathrm{GL}_m(K)$。由于该伪表示是 $\sigma$-不变的，这样的 $\rho$ 可延拓为某个有限扩张 $K'/K$ 上的连续表示 $W_{F,S}\to\mathrm{GL}_m(K')$。于是由 de Jong 猜想，$\pi_1(\overline Y)$ 的像有限。因此 $B^{\Theta}\to\kappa'[[t]]$ 的像是 $\kappa'$。这就说明 $B^{\Theta}/\ell$ 是阿廷环。
\end{proof}

现在证明定理 3.29。
\begin{proof}
我们使用 Artin–Lurie 可表示性定理 [57, 7.5.1]。先验证 $\mathcal R^{sc}_{\pi_1(\overline Y),{}^{c}G}$ 满足 loc. cit. 的条件 (1)–(5)。即：$\mathcal R^{sc}_{\pi_1(\overline Y),{}^{c}G}$ 是 $0$-截断的，故条件 (2) 成立；由命题 2.59，条件 (1)、(4)、(5) 成立。我们断言 $\mathcal R^{sc}_{\pi_1(\overline Y),{}^{c}G}$ 满足 fppf 下降，故条件 (3) 也成立。事实上，由于它是幂零完备的，只需证明
\[
\mathcal R^{sc}_{\pi_1(\overline Y),{}^{c}G}(A)\to\varprojlim_{\Delta}\mathcal R^{sc}_{\pi_1(\overline Y),{}^{c}G}(B^{\bullet})
\]
是同构，其中 $B^{\bullet}:\Delta\to{}_{\le m}\mathbf{CAlg}_{\Z_{\ell}}$ 是 $m$-截断动态（animated）$\Z_{\ell}$-代数之间忠实平坦映射 $A\to B$ 的 Čech 神经。此时可以把 $\Delta$ 上的极限换成由对象
\ysrc{86}
$[0],\dots,[m+1]$ 组成的 $\Delta_{\le m+1}\subset\Delta$ 上的有限极限。由于 $\mathcal R^{sc}_{\pi_1(\overline Y),{}^{c}G}$ 保持有限极限，断言成立。现在容易看出 $\Loc_{{}^{c}G,F,S}\to\mathcal R^{sc}_{\pi_1(\overline Y),{}^{c}G}/\hat G$ 相对可表示，故 $\Loc_{{}^{c}G,F,S}$ 也满足 [57, 7.5.1] 的条件 (1)–(5)。

{\sloppy
再由命题 2.59，在点 $\rho:W_{F,S}\to{}^cG(A)$ 处，$\Loc_{{}^{c}G,F,S}$ 的切空间是连续上同调 $C^{*}_{cts}(W_{F,S},\mathrm{Ad}^{0}_{\rho})[1]$，其中 $A$ 是经典 $\Z_{\ell}$-代数，$\mathrm{Ad}^{0}$ 是 ${}^cG$ 在 $\hat G$ 的李代数上的伴随表示。回忆：对 $\pi_1(\overline Y)$ 在有限 $\Z_{\ell}$-模 $V$ 上的连续表示，连续群上同调 $C^{*}_{cts}(\pi_1(\overline Y),V)$ 同构于 $V$ 的 étale 上同调（把 $V$ 视为仿射簇 $\overline Y$ 上的局部系统）。由引理 2.57 与 (2.40) 可知 $C^{*}_{cts}(\overline Y,\mathrm{Ad}^{0}_{\rho})$ 集中在次数 $[-1,0]$，且当 $A$ 在 $\Z_{\ell}$ 上有限生成时各上同调群是有限 $A$-模。于是 Hochschild–Serre 谱序列蕴含 $C^{*}_{cts}(W_{F,S},\mathrm{Ad}^{0}_{\rho})[1]$ 集中在次数 $[-1,1]$，且当 $A$ 在 $\Z_{\ell}$ 上有限生成时每个次数上都是有限 $A$-模。这验证了 [57, 7.5.1] 的条件 (7)。此外，它还表明若 $\Loc_{{}^{c}G,F,S}$ 可表示，则它拟光滑。\par}

剩下验证条件 (6)。我们证明：对经典诺特完备 $\Z_{\ell}$-代数 $(A,\mathfrak m)$，其剩余域 $\kappa$ 在 $\mathbb F_{\ell}$ 上或在 $\Q_{\ell}$ 上有限，映射
\[
\Loc^{\square}_{{}^{c}G,F,S}(A)\to\varprojlim_{i}\Loc^{\square}_{{}^{c}G,F,S}(A/\mathfrak m^{i})
\]
是同构。取忠实表示 ${}^cG\to\mathrm{GL}_m$ 后，问题归约为证明
\begin{equation}
\mathcal R^{sc}_{W_{F,S},\mathrm{GL}_m}(A)\to\varprojlim_{i}\mathcal R^{sc}_{W_{F,S},\mathrm{GL}_m}(A/\mathfrak m^{i}) \tag{3.26}
\end{equation}
是同构。设 $\{\rho_i\}$ 是一族相容的表示 $\rho_i:W_{F,S}\to\mathrm{GL}_m(A/\mathfrak m^{i})$，给出 (3.26) 右端的一个元素。注意由于 $A/\mathfrak m^{i}$ 在 $\Z_{\ell}$ 上或在 $\Q_{\ell}$ 上有限，每个 $\rho_i$ 就是通常意义下的连续表示（见注 2.54）。忘掉拓扑并取逆向极限，得到表示 $\rho:W_{F,S}\to\mathrm{GL}_m(A)$。需要证明它强连续。由引理 2.57，只需证明对每个 $v\in A^{m}$，$\rho(\pi_1(\overline Y))v$ 包含在一个有限 $\Z_{\ell}$-模中。

设 $B$ 是 $A$ 的 $\Z_{\ell}$-子代数，由 $\chi_j(\rho(\gamma))$（$\gamma\in\pi_1(\overline Y)$）生成，其中 $\chi_i\in\Z[\mathrm{GL}_m]^{\mathrm{GL}_m}$ 如前是 $\mathrm{GL}_m$ 的第 $i$ 个楔表示的特征标。则对每个 $\gamma\in\pi_1(\overline Y)$，特征多项式 $\mathrm{Char}(\rho(\gamma),t)\in B[t]$。把 $\pi_1(\overline Y)$ 在 $A^{m}$ 上的作用延拓为群环 $B\pi_1(\overline Y)$ 的作用。注意 $r=\sum b_j\gamma_j\in B\pi_1(\overline Y)$ 的特征多项式也属于 $B[t]$。由于每个 $\rho_i$ 连续（通常意义下），该作用延拓为完备群环 $B\pi_1(\overline Y)^{\wedge}$ 的作用，继而通过商 $B\pi_1(\overline Y)^{\wedge}/I$ 分解，其中 $I$ 是由 $\mathrm{Char}(\rho(r),r)$（$r\in B\pi_1(\overline Y)^{\wedge}$）生成的理想。由于 $\pi_1(\overline Y)$ 满足 Mazur 条件 $\Phi_{\ell}$，由 [75, 3.6]，$B\pi_1(\overline Y)^{\wedge}/I$ 在 $B$ 上有限。我们断言 $B$ 在 $\Z_{\ell}$ 上有限，这将完成 (3.26) 是同构的证明。

考虑 $\rho_0:W_{F,S}\to\mathrm{GL}_m(A/\mathfrak m)=\mathrm{GL}_m(\kappa)$。若 $\kappa$ 是有限域，令 $\bar\rho=\rho_0|_{\pi_1(\overline Y)}$。若 $\kappa=E$ 是特征零的，则共轭后可以假定 $\rho_0|_{\pi_1(\overline Y)}$ 来自一个 $\mathcal O_E$-表示。令 $\bar\rho:\pi_1(\overline Y)\to\mathrm{GL}_m(\kappa_E)$ 是 $\rho_0|_{\pi_1(\overline Y)}$ 的剩余表示。有通常的（经典）加框架形变环 $R^{\square}_{\bar\rho}$。表示 $\rho_0|_{\pi_1(\overline Y)}$ 给出 $R^{\square}_{\bar\rho}$ 的一个点，而
\ysrc{87}
$R^{\square}_{\bar\rho}$ 在该点的形式完备化预表示 $\rho_0$ 的经典加框架形变（视为函子 $\mathbf{Art}_{\Z_{\ell},\kappa}\to\mathbf{Sets}$）。（若 $\kappa=E$，见 [44, 2.3.5]。）于是 $\rho|_{\pi_1(\overline Y)}:\pi_1(\overline Y)\to\mathrm{GL}_m(A)$ 给出映射 $R^{\square}_{\bar\rho}\to A$。设 $\Theta$ 是与 $\bar\rho$ 相伴的伪表示。则有引理 3.30 中的 $A^{\Theta}$，而 $B$ 恰是 $A^{\Theta}$ 在自然映射 $A^{\Theta}\to R^{\square}_{\bar\rho}\to A$ 下的像，该映射通过 $A^{\Theta}/(1-\sigma)A^{\Theta}\to A$ 分解。因此由引理 3.30，$B$ 在 $\Z_{\ell}$ 上有限。

我们已证明 $\Loc_{{}^{c}G,F,S}$ 的可表示性。由引理 2.60 得到分解 (3.25)。还需看出 $\Loc^{\Theta}_{{}^{c}G,F,S}$ 拟紧。事实上我们要证明相应的加框架版本 $\Loc^{\Theta,\square}_{{}^{c}G,F,S}$ 由 $\Z_{\ell}$ 上有限型的仿射概形表示。可以归约到 $\mathrm{GL}_m$ 情形。有引理 3.30 中的环 $B^{\Theta}$，以及如前所述的有限（结合）$\Z_{\ell}$-代数 $B^{\Theta}\pi_1(\overline Y)^{\wedge}/I$。把 Frobenius $\sigma$ 提升为 $W_{F,S}$ 中的元素，于是 $\sigma$ 作用在 $B^{\Theta}\pi_1(\overline Y)^{\wedge}/I$ 上，可以构造扭 $\Z_{\ell}$-代数 $B^{\Theta}\pi_1(\overline Y)^{\wedge}/I[\sigma]$。此时 ${}^{cl}\Loc^{\Theta,\square}_{{}^{c}G,F,S}$ 恰是有限生成结合 $\Z_{\ell}$-代数 $B^{\Theta}\pi_1(\overline Y)^{\wedge}/I[\sigma]$ 的加框架 $m$ 维表示构成的模空间，因此由 $\Z_{\ell}$ 上有限型的仿射概形表示。
\end{proof}

\begin{rem}{3.31}
可以把分解 (3.25) 视为局部情形中 mod $\ell$ 惯性型（注 3.3）的整体类比。显然 $\Loc^{\Theta}_{{}^{c}G,F,S}$ 非空当且仅当 $\Theta$ 在 Frobenius $\sigma$ 的作用下不动。
\end{rem}

\begin{rem}{3.32}
如局部情形那样，人们可能期望叠 $\Loc_{{}^{c}G,F,S}$ 是经典的。如注 2.34 所述，$\Loc^{\Theta}_{{}^{c}G,F,S}$ 经典当且仅当 $\dim\Loc^{\Theta}_{{}^{c}G,F,S}=0$。遗憾的是，这并不总成立。

考虑 $G=\mathrm{PGL}_2$（从而 ${}^cG=\mathrm{GL}_2$）的情形，设 $\Theta$ 是与 $\pi_1(\overline Y)$ 的平凡表示相应的伪表示。则 $\Loc^{\Theta}_{{}^{c}G,F,S,1}$ 由那些 $\rho:W_{F,S}\to\mathrm{GL}_2$ 组成，使得 $\rho|_{\pi_1(\overline Y)}$ 是平凡特征标的自扩张。注意 $\pi_1(\overline Y)$ 的平凡特征标的自扩张构成 $H^{1}(\overline Y,\mathbb F_{\ell})$ 族。由此可见，若某个 Frobenius 特征值在 $H^{1}(\overline Y,\mathbb F_{\ell})$ 上的重数大于一，则 $\dim\Loc^{\Theta,\square}_{{}^{c}G,F,S,1}>\dim\hat G_{\mathbb F_{\ell}}$，且 $\Loc^{\Theta}_{{}^{c}G,F,S,1}$ 非经典。
\end{rem}

有时考虑具有固定「行列式」的 $\Loc_{{}^{c}G,F,S}$ 的子叠是方便的。更确切地说，设 $Z^{\circ}_G$ 是 $G$ 的连通中心。则 ${}^c(Z^{\circ}_G)=\hat G_{ab}\rtimes(\mathbb G_m\times\Gamma_{\wt F/F})$，其中 $\hat G_{ab}$ 是 $\hat G$ 的交换化。存在自然态射 $\pi_{ab}:\Loc_{{}^{c}G,F,S}\to\Loc_{{}^{c}(Z^{\circ}_G),F,S}$。给定经典 $\Z_{\ell}$-代数 $A$ 与强连续表示 $\lambda:W_{F,S}\to{}^c(Z^{\circ}_G)(A)$（满足 $d\circ\lambda=\chi$），它对应于 $\Loc_{{}^{c}(Z^{\circ}_G),F,S}$ 的一个 $A$-点；令
\[
\Loc^{\lambda}_{{}^{c}G,F,S,A}=\Loc_{{}^{c}G,F,S}\times_{\Loc_{{}^{c}(Z^{\circ}_G),F,S}}\Spec A
\]
记 $\pi_{ab}$ 沿 $\lambda$ 的底变，它是 $A$ 上的代数叠，分类那些满足 $\pi_{ab}\circ\rho=\lambda$ 的表示 $\rho$。它在 $\rho$ 处的切空间由 $C^{*}_{cts}(W_{F,S},\mathrm{Ad}^{00})$ 给出，其中 $\mathrm{Ad}^{00}$ 是 ${}^cG$ 在 $\hat G$ 的导群李代数上的伴随表示。特别地，$\Loc^{\lambda}_{{}^{c}G,F,S,A}$ 在 $A$ 上拟光滑。
\begin{exam}{3.33}
椭圆 Langlands 参数是连续半单表示 $\rho:W_{F,S}\to{}^cG(\overline\Q_{\ell})$（满足 $d\circ\rho=\chi$），使得 $\mathfrak S_{\rho}:=S_{\rho}/(Z_{\hat G})^{\Gamma_F}$ 有限，其中 $S_{\rho}$ 是 $\rho$ 在 $\hat G$ 对 ${}^cG$ 的共轭作用下的稳定化子，
\ysrc{88}
而 $Z_{\hat G}$ 是 $\hat G$ 的中心，$W_F$ 在其上有作用。由 [50, 4.1]，椭圆 Langlands 参数 $\rho$ 在 $\Loc^{\lambda}_{{}^{c}G,F,S,\overline\Q_{\ell}}$ 中给出孤立光滑点，其中 $\lambda=\pi_{ab}\circ\rho$。更确切地说，每个椭圆 $\rho$ 给出开闭嵌入 $(\Spec\overline\Q_{\ell})/\mathfrak S_{\rho}\hookrightarrow\Loc^{\lambda}_{{}^{c}G,F,S,\overline\Q_{\ell}}$。
\end{exam}

把 $W_{F_v}\to W_F$ 的嵌入模去共轭，诱导出良定义的态射
\begin{equation}
\mathrm{res}:\Loc_{{}^{c}G,F,S}\to\prod_{v\in S}\Loc_v\times\prod_{w\notin S}\Loc^{\mathrm{ur}}_{w}. \tag{3.27}
\end{equation}
\begin{lem}{3.34}
下图中的交换方块是笛卡尔（拉回）的：
\begin{equation}
\begin{tikzcd}[column sep=large,row sep=large]
\Loc_{{}^{c}G,F,S} \arrow[r] \arrow[d] & \prod_{v\in S}\Loc_v\times\Loc^{\mathrm{ur}}_{w_0} \arrow[r] \arrow[d] & \prod_{v\in S}\Loc_v \\
\Loc_{{}^{c}G,F,S\cup\{w_0\}} \arrow[r] & \prod_{v\in S}\Loc_v\times\Loc_{w_0} \arrow[r] & \prod_{v\in S}\Loc_v
\end{tikzcd} \tag{3.28}
\end{equation}
\end{lem}
\begin{proof}
由幂零完备性，只需证明该图在取值于 $m$-截断动态 $\Z_{\ell}$-代数 $A$ 时是笛卡尔方块。当 $A$ 是经典代数时这是显然的。然后用 Postnikov 塔并按命题 2.42 的方式论证，就归约为比较切空间，这并不困难。细节留给读者。（在密切相关的情形下可参见 [35, §8] 的论证。）
\end{proof}

对每个素位 $v\in S$，选取一个在 $F_v$ 上 Galois 的有限扩张 $L_v/F^{t}_v\wt F_v$。令
\[
\Loc_{{}^{c}G,F,\{L_v\}}:=\Loc_{{}^{c}G,F,S}\times_{\prod_{v\in S}\Loc_v}\prod_{v\in S}\Loc_{{}^{c}G_v,L_v/F_v}.
\]
由于 $\Loc_{{}^{c}G_v,L_v/F_v}$ 在 $\Loc_v$ 中开闭，叠 $\Loc_{{}^{c}G,F,\{L_v\}}$ 也在 $\Loc_{{}^{c}G,F,S}$ 中开闭。特别地，若 $G$ 在 $F$ 上驯顺分歧，则有驯顺叠 $\Loc^{\mathrm{tame}}_{{}^{c}G,F,S}:=\Loc_{{}^{c}G,F,\{F^{t}_v\}}$。
\begin{prop}{3.35}
叠 $\Loc_{{}^{c}G,F,\{L_v\}}$ 拟紧，且 $\Loc_{{}^{c}G,F,S}=\cup_{\{L_v\}}\Loc_{{}^{c}G,F,\{L_v\}}$。
\end{prop}
\begin{proof}
可以忽略导出结构。把 $W_{F,S}$（相应地 $\Gamma_{F,S}$）由子群族 $\{\Gamma_{L_v},v\in S\}$ 的（共轭类）所生成的闭正规子群取商，记作 $W_{F,\{L_v\}}$（相应地 $\Gamma_{F,\{L_v\}}$）。对固定的忠实表示 ${}^cG\to\mathrm{GL}_m$，诱导态射 ${}^{cl}\mathcal R^{sc}_{W_{F,\{L_v\}},{}^{c}G}\to{}^{cl}\mathcal R^{sc}_{W_{F,\{L_v\}},\mathrm{GL}_m}={}^{cl}\Loc^{\square}_{{}^{c}(\mathrm{GL}_m),F,\{L_v\}}$ 是闭嵌入。因此只需对 $G=\mathrm{GL}_m$ 证明该命题。

现在分解 (3.25) 给出分解 $\Loc_{{}^{c}(\mathrm{GL}_m),F,\{L_v\}}=\sqcup_{\Theta}\Loc^{\Theta}_{{}^{c}(\mathrm{GL}_m),F,\{L_v\}}$，故只需证明分解中只出现有限多个这样的 $\Theta$。每个这样的 $\Theta$ 给出连续半单表示 $\bar\rho:\Gamma_{F,\{L_v\}}\to\mathrm{GL}_m(\overline{\mathbb F}_{\ell})$；对 $\bar\rho$ 的每个不可约因子应用 [22, 3.5]，可知它可提升为特征零中具有有限行列式的半单表示 $\rho$。（注意由于 $S$ 非空，[22, 3.5] 的假设 (iii) 是不必要的。）由 L. Lafforgue [48] 证明的函数域上 $\mathrm{GL}_m$ 的整体 Langlands 对应，这样的 $\rho$ 在共轭意义下只有有限多个。
\end{proof}

\ysrc{89}
\begin{rem}{3.36}
注意在 $\Loc_{{}^{c}G,F,S}$ 的定义中我们总要求 $S$ 是非空的有限集（以保证连续群上同调与 étale 上同调一致）。这先验地排除了处处非分歧的 Langlands 参数的叠。不过引理 3.34 使我们能如下恢复这种情形。假定 $\Gamma_F$ 在 $\hat G$ 上的作用通过非分歧伽罗瓦群分解，即分解到以 $F$ 为分式域的 $\mathbb F_q$ 上光滑射影曲线 $X$ 的 étale 基本群 $\pi_1(X)$。设 $S=\{v\}$ 是 $X$ 的一个素位。则定义
\[
\Loc_{{}^{c}G,X}:=\Loc_{{}^{c}G,F,\emptyset}:=\Loc_{{}^{c}G,F,\{v\}}\times_{\Loc_v}\Loc^{\mathrm{ur}}_v=\Loc^{\mathrm{tame}}_{{}^{c}G,F,\{v\}}\times_{\Loc^{\mathrm{tame}}_v}\Loc^{\mathrm{ur}}_v.
\]
它与 $v$ 的选取无关。例如若 $X=\PP^1$，则
\begin{equation}
\Loc_{{}^{c}G,\PP^1}=\Loc_{{}^{c}G,F,\{\infty\}}\times_{\Loc_{\infty}}\Loc^{\mathrm{ur}}_{\infty}\cong\Loc^{\mathrm{ur}}_{0}\times_{\Loc^{\mathrm{tame}}_{{}^{c}G,F,\{0,\infty\}}}\Loc^{\mathrm{ur}}_{\infty}. \tag{3.29}
\end{equation}
由命题 3.35，显然 $\Loc_{{}^{c}G,X}$ 拟紧。当 $S=\emptyset$ 时椭圆参数的概念仍有意义，它们仍在相应的 $\Loc^{\lambda}_{{}^{c}G,X,\overline\Q_{\ell}}$ 中给出孤立光滑点。

在本小节末尾，简要提及 $F$ 是数域的情形，这是与 Emerton 的进行中合作 [26]。此时仍有 $\chi:\Gamma_{F,S}\to\Z_{\ell}^{\times}\times\Gamma_{\wt F/F}$，其中第一个分量是分圆特征标的逆。把它视为 $\mathcal R_{\Gamma_{F,S},\mathbb G_m\times\Gamma_{\wt F/F}}$ 的一个 $\Spf\Z_{\ell}$-点。于是与 (3.23) 类似，令
\[
\Loc^{\wedge,\square}_{{}^{c}G,F,S}:=\mathcal R_{\Gamma_{F,S},{}^{c}G}\times_{\mathcal R_{\Gamma_{F,S},\mathbb G_m\times\Gamma_{\wt F/F}}}\{\chi\},\qquad
\Loc^{\wedge}_{{}^{c}G,F,S}=\Loc^{\wedge,\square}_{{}^{c}G,F,S}/\hat G^{\wedge}_{\ell}.
\]
仍有 $\Loc^{\wedge}_{{}^{c}G,F,S}=\varinjlim_{r}\Loc_{{}^{c}G,F,S,r}$，其中 $\Loc_{{}^{c}G,F,S,r}$ 是 $\Loc_{{}^{c}G,F,S}$ 到 $\Z/\ell^{r}$ 的限制。不过对数域情形情况更复杂。首先连 $\Loc_{{}^{c}G,F,S,1}$ 一般也不是代数叠，而只是 ind-叠。此外，数域情形下我们不会用定义 2.56 去定义 $\Z_{\ell}$ 上的叠，因为这样的对象可能并不合理。我们转而考虑整体到局部的态射
\[
\mathrm{res}:\Loc^{\wedge}_{{}^{c}G,F,S}\to\prod_{v\in S}\Loc^{\wedge}_v,
\]
其中若 $v$ 不在 $\ell$ 之上，$\Loc^{\wedge}_v$ 如 (3.2) 所定义；若 $v$ 在 $\ell$ 之上（比如 ${}^cG=\mathrm{GL}_n$），则取 [23] 中的叠。在 [26] 中我们将证明：在 de Jong 猜想的一个类比之下，这个态射可表示。这样的结论对许多应用应当足够，例如给出志村簇上同调的一个猜想性公式。利用这个态射，可以在 $v\mid\ell$ 处施加 $\ell$-进 Hodge 理论条件（例如具有某些固定 Hodge–Tate 权的晶体性），从 $\Loc^{\wedge}_{{}^{c}G,F,S}$ 中截出闭子叠，它们将是 $\ell$-进形式叠。这些子叠随后可能容许延拓为 $\Z_{\ell}$ 上的代数叠，那将是数域情形下 $\Loc_{{}^{c}G,F,S}$ 的正确类比。
\end{rem}

\ysrc{90}
\section{Langlands 参数叠上的凝聚层}

本节利用 Langlands 参数的叠来表述局部与整体 Langlands 对应中的一些猜想。我们也综述一些已知结果，它们为这些猜想提供证据。本节中 $\Lambda$ 表示诺特交换环。

本节出现的许多范畴将是 $\Lambda$-线性稳定 $\infty$-范畴（见 [55, Chap 1]）。对这样的范畴 $\mathcal C$ 中的两个对象 $x_1,x_2$，它们的（导出）hom 空间自然地是 $\Lambda$-模，记作 $\mathrm{Hom}_{\mathcal C}(x_1,x_2)$（若 $\mathcal C$ 由上下文清楚，也简记作 $\mathrm{Hom}(x_1,x_2)$）。原来的映射空间 $\Map_{\mathcal C}(x_1,x_2)$ 与 $\tau^{\le 0}\mathrm{Hom}_{\mathcal C}(x_1,x_2)$ 等同。按记号滥用，我们用 $\End(x)$ 表示 $\mathrm{Hom}(x,x)$，它是 $\mathbf{Alg}(\mathbf{Mod}_{\Lambda})$ 中的对象，即 $E_1$-代数。（注意在例 2.2 中我们用同一记号表示 $x$ 的自同态幺半群。希望该记号的具体含义可由上下文看清。）

\subsection{$G(F)$ 的表示范畴}

设 $F$ 是非阿基米德局部域，$\mathcal O_F$ 是其整数环，$\kappa_F$ 其剩余域。还固定一个单值化元 $\varpi_F\in\mathcal O_F$。令 $q=\sharp\kappa_F=p^{r}$。设 $G$ 是 $F$ 上的连通约化群。令 $\Rep(G(F),\Lambda)^{\heartsuit}$ 记 $G(F)$ 在 $\Lambda$-模上的光滑表示的阿贝尔范畴。它是 Grothendieck 阿贝尔范畴（下文给出一组生成元）。对闭子群 $K\subset G(F)$，类似地有 $\Rep(K,\Lambda)^{\heartsuit}$。我们总用 $\mathbf 1$ 表示平凡表示。令
\[
c\text{-ind}^{G(F)}_{K}:\Rep(K,\Lambda)^{\heartsuit}\to\Rep(G(F),\Lambda)^{\heartsuit}
\]
记通常的紧诱导函子，并记
\[
\delta_K:=c\text{-ind}^{G(F)}_{K}\mathbf 1\cong C^{\infty}_c(G(F)/K,\Lambda),
\]
它是 $G(F)/K$ 上紧支集的 $\Lambda$-值局部常值函数空间，$G(F)$ 通过左平移作用在其上。

若 $K$ 开，则 $c\text{-ind}^{G(F)}_{K}$ 是遗忘函子的左伴随。由光滑表示的定义，族 $\{\delta_K\}$（$K$ 取遍开子群）构成 $\Rep(G(F),\Lambda)^{\heartsuit}$ 的一组生成元。称 $G(F)$ 的开紧子群 $K$ 为 $\Lambda$-容许的（若 $\Lambda$ 由上下文清楚，也简称容许），如果 $K$ 的任意开子群的指数在 $\Lambda$ 中可逆。注意若 $p$ 在 $\Lambda$ 中可逆，则 $\Lambda$-容许开紧子群总存在。例如 $G(F)$ 的 Iwahori 子群（有时也称 pro-p Iwahori 子群）的 pro-$p$ Sylow 子群 $I(1)$ 是 $\Lambda$-容许的。另一方面，每个开紧子群都是 $\Q$-容许的。若 $K$ 是 $\Lambda$-容许的，则 $\delta_K$ 是 $\Rep(G(F),\Lambda)^{\heartsuit}$ 中的投射对象。

其次，令 $\Rep(G(F),\Lambda)$ 记 $\Rep(G(F),\Lambda)^{\heartsuit}$ 的（无界）$\infty$-导出范畴（[55, 1.3.5]）。该范畴的行为在 $p$ 是否于 $\Lambda$ 中可逆时相当不同。为我们的目的，本节通篇假定 $p$ 在 $\Lambda$ 中可逆。此时 $c\text{-ind}^{G(F)}_{K}$ 是 $t$-正合函子。若 $K$ 是 $\Lambda$-容许开紧子群，则 $\delta_K$ 是 $\Rep(G(F),\Lambda)$ 中的紧对象。由此 $\Rep(G(F),\Lambda)$ 是紧生成的，一组生成元由 $\{\delta_K\}$（$K$ 取遍 $\Lambda$-容许开紧子群）给出。
\begin{rem}{4.1}
若 $F$ 是特征零的，而 $\Lambda$ 是特征 $p$ 的域（这不是我们考虑的情形），则 $\delta_{I(1)}$ 本身就是 $\Rep(G(F),\Lambda)$ 的紧生成元（见 [64]）。
\end{rem}

一般来说，若开紧子群 $K$ 不是 $\Lambda$-容许的，则 $\delta_K$ 在 $\Rep(G(F),\Lambda)$ 中可能不紧。
\begin{exam}{4.2}
若 $G=\mathbb G_m$，$K=\mathcal O_F^{\times}$，$\Lambda=\mathbb F_{\ell}$，其中 $\ell$ 是整除 $q-1$ 的素数，则 $\delta_K\simeq C_c(\Z,\mathbb F_{\ell})$ 在 $\Rep(F^{\times},\mathbb F_{\ell})$ 中不紧。
\end{exam}
